Let \(p_c=2n/(n-1)\). By \ref{ms:thm-1-2} and \ref{ms:lemma-9-2}, local spherical extension \(L^{p_c}(S^{n-1})\to L^{p_c}(B(0,R))\) has norm \(O([\log(2+R)]^{2/p_c})\), hence \(O_\epsilon(R^\epsilon)\) for each \(\epsilon>0\). Duality gives the same loss for restriction \(f\mapsto\widehat f|_{S^{n-1}}\) from \(L^{p_c'}(\mathbb R^n)\) for \(f\) supported in \(B(0,R)\), with values in \(L^{p_c'}(S^{n-1})\). Here \(\widehat f(\xi)=\int e^{-ix\cdot\xi}f(x)\,dx\). Apply the epsilon-removal theorem of \cite{r-75}, Theorem 1.2. Its consequence used here is that local \(L^u\to L^u\) restriction on the sphere at every loss \(O_\epsilon(R^\epsilon)\), for \(1<u<2\), yields global \(L^v(\mathbb R^n)\to L^u(S^{n-1})\) bounds for each \(1<v<u\). Since \(p_c'\in(1,2)\), this gives global diagonal \(L^v\to L^v\) restriction by finite sphere measure, and therefore by duality spherical extension \(L^s\to L^s\) for every \(p_c<s<\infty\).

If \(a\ge b\), take \(s=b\) and use finite measure of the sphere. If \(a<b\), set \(s=1+b/a'\). The range hypotheses give \(p_c<s<b\). Interpolate the diagonal estimate and the defining integral bound \(E:L^1\to L^\infty\) with weight \(\theta=s/b\) on the diagonal estimate. The interpolated input and output reciprocals are \(1-\theta+\theta/s=1/a\) and \(\theta/s=1/b\), so the spherical estimate holds.

For finite \(b\) the hypotheses imply \(b>2\) and \(b>a'\), hence \(a>b'\). The spherical bound passes to the cone by \ref{ms:cone-slicing}. For \(b=\infty\) each defining integral is bounded by the \(L^1\) norm, which is bounded by the \(L^a\) norm on a compact surface.
